\documentclass[a4paper,10pt]{article}
\usepackage[dutch]{babel}
\usepackage[T1]{fontenc}
\usepackage{lmodern}
\usepackage[utf8]{inputenc}
\usepackage[margin=1.5cm]{geometry}
\usepackage{amsmath}
\usepackage{hyperref}
% Avoid rerunfilecheck warning about changed .out outlines
% (see message: use package `bookmark')
\usepackage{bookmark}
\usepackage{enumitem}
\usepackage{booktabs}
\usepackage{multicol}
\usepackage{xcolor}
\usepackage{xspace}
\usepackage{mdframed}
\usepackage{titlesec}
\usepackage[nofontexpansion]{microtype}

\titleformat{\section}{\large\bfseries\color{blue}}{}{0em}{}[\titlerule]
\titleformat{\subsection}{\normalsize\bfseries\color{teal}}{}{0em}{}
\titleformat{\subsubsection}{\small\bfseries\color{orange!70!black}}{}{0em}{}

\setlength{\parskip}{3pt}
\setlength{\parindent}{0pt}

\newmdenv[backgroundcolor=yellow!20, linecolor=orange, linewidth=1pt, innertopmargin=4pt, innerbottommargin=4pt]{exambox}
\newmdenv[backgroundcolor=blue!10, linecolor=blue, linewidth=1pt, innertopmargin=4pt, innerbottommargin=4pt]{infobox}
\newmdenv[backgroundcolor=green!10, linecolor=green!60!black, linewidth=1pt, innertopmargin=4pt, innerbottommargin=4pt]{answerbox}

\begin{document}

\begin{center}
{\LARGE\bfseries RELIGIONS EXAMEN SPIEKBOEK}\\[4pt]
{\large A07M9a -- Living in the Age of Humans}\\[2pt]
{\normalsize Examinator: Depoortere Frederiek $\cdot$ 120 min $\cdot$ Open boek}\\[2pt]
{\small \textbf{REGELS 2025:} Annotaties in boek toegelaten + kleine post-its. GEEN losse bladen. Nederlands mag.}
\end{center}

\begin{exambox}
\textbf{EXAMENSTRUCTUUR (2025, 17/20):}
\begin{itemize}[noitemsep,topsep=2pt]
  \item \textbf{Vraag 1 (4/20):} 4 korte definities (elk $\sim$3 regels). Antwoord staat in het boek.
  \item \textbf{Vraag 2 (10/20):} Synthese/reflectie op 1 sleutelvraag. Max 2 pagina's. Verwijzen naar het boek!
  \item \textbf{Vraag 3 (3/20):} Persoonlijke reflectie op het leerproces. Minstens 1 pagina.
\end{itemize}
\textbf{STRATEGIE:} Schrijf de uitgewerkte antwoorden op de sleutelvragen op de relevante pagina's van het boek. Schrijf definities bij de betreffende secties. Markeer pagina's met post-its.
\end{exambox}

\vspace{4pt}

%=========================================================
\section{WAT TE SCHRIJVEN EN WAAR (PER HOOFDSTUK)}
%=========================================================

\begin{infobox}
\textbf{STRUCTUUR VAN HET BOEK:}
\begin{itemize}[noitemsep]
  \item \textbf{Hfst 1--2:} De Crisis (ecologische crisis als spirituele crisis)
  \item \textbf{Hfst 3--4:} Wetenschap en Religie (NOMA, conflict, scheiding, integratie -- Barbour)
  \item \textbf{Hfst 5:} De spirituele uitdaging / Technologie en spiritualiteit
  \item \textbf{Hfst 6:} Vergroening van de wereldreligies (Jodendom, Islam, Christendom, Oosterse religies)
  \item \textbf{Hfst 7:} Aardse systeemwetenschap als spirituele bron (Gaia)
  \item \textbf{Hfst 8:} Environmentalisme als religie
\end{itemize}
\end{infobox}

\vspace{3pt}

\subsection{VOORPAGINA / INLEIDING (schrijf hier de definities van kernbegrippen)}

\textbf{Op de voorste lege pagina's of het schutblad schrijven:}

\textbf{SPIRITUALITEIT} (p.22--28): The personal search for meaning, purpose, and connectedness to a larger whole. Three key attributes: (1) \textit{connectedness} -- sense of belonging to something larger; (2) \textit{self-transcendence} -- going beyond oneself; (3) \textit{meaning} -- purpose in life. NOT necessarily tied to belief in God.

\textbf{RELIGIE} (p.22--28): A set of beliefs and behaviours connecting humans to a perceived higher order or unseen dimension. Religion expands on spirituality by implying a higher power, a moral system, and ritual practice. Religion typically involves a community.

\textbf{IDEOLOGISCH PROFIEL}: The set of beliefs and values that shape a person's worldview; presuppositions that underlie how one interprets reality.

\textbf{ANTHROPOCEEN}: The current geological epoch in which human activity is the dominant influence on climate and environment; humans as a ``geological superpower.''

\textbf{COSMOS}: The universe seen as a well-ordered whole.

\vspace{3pt}

\subsection{HOOFDSTUK 1--2: De Ecologische Crisis als Spirituele Crisis}

\textbf{Schrijf bij het begin van hfst 1:} $\to$ \textbf{SLEUTELVRAAG 1: "Is de ecologische crisis een spirituele crisis?"} Zie uitgewerkt antwoord verderop.

\textbf{Kernidee (schrijf bij sectie 1.x):} The ecological crisis is not merely a technical or political problem but fundamentally a \textit{spiritual crisis}: a crisis of meaning, values and our self-understanding as human beings. The root cause is a cultural and philosophical shift in how we understand our place in the cosmos.

\textbf{White Thesis} (Lynn White Jr., 1967, ``The Historical Roots of Our Ecologic Crisis''): Western Christianity -- specifically \textit{Latin} (Western) Christianity -- is the most anthropocentric religion in history. It desacralized nature: transformed nature from an animated world of protective spirits into ``lifeless and soulless objects'' designed to serve human ends. \textit{Voluntaristic} (orthopraxis, right conduct over contemplation) vs.\ \textit{Intellectualistic} (Greek/Eastern Christianity: orthodoxy, contemplation, nature as symbol of the divine).

\textbf{Vier reacties van Abrahamitische religies op White:}
\begin{itemize}[noitemsep]
  \item \textit{Hostile Rejection}: environmentalism as ``new paganism''
  \item \textit{Indifference}: ecology not a religious concern
  \item \textit{Apologetics}: tradition misunderstood, supports nature
  \item \textit{Confession of Guilt}: recognising historical failings, seeking reform
\end{itemize}

\vspace{3pt}

\subsection{HOOFDSTUK 2 (§2.3.1, p.51): Relatie Religie en Spiritualiteit}

\textbf{Schrijf op p.51:} Spirituality is the broader concept; religion is a specific organized form of spirituality. Religion can be a source of spirituality, but spirituality does not require religion. In the Anthropocene, natural science is increasingly a \textit{new source of spirituality} (awe, connectedness, meaning through scientific discovery).

\textbf{§2.3.2 Spirituality and Natural Science:} Natural science as source of spirituality through: (1) connectedness -- the cosmos as one, Pale Blue Dot/Carl Sagan; (2) self-transcendence -- science humbles us, we are ``tiny twigs''; (3) meaning -- the universe's story gives humans a narrative.\\
\textbf{Carl Sagan quote (examenvraag aug 2025!):} ``Science is not only compatible with spirituality; it is a profound source of spirituality.'' $\to$ Because science generates awe (connectedness to cosmos), humility (self-transcendence), and narrative meaning.

\textbf{§2.3.3 Technologie en Spiritualiteit} (p.65 Miss Case):
Miss Case: technology makes us more human by connecting us and giving access to vast data $\to$ if we find spirituality all around us, technology is a form of spirituality. BUT: danger of the \textit{technocratic paradigm} (Francis).

\vspace{3pt}

\subsection{HOOFDSTUK 3 (p.76--80): Wetenschap en Religie -- Barbour's Modellen}

\textbf{Schrijf op de eerste pagina van hfst 3 (groot en duidelijk):}

\textbf{BARBOUR'S 4 MODELLEN} (Ian Barbour, \textit{Issues in Science and Religion}, 1966):
\begin{enumerate}[noitemsep]
  \item \textbf{Conflict}: Science and religion are ``sworn enemies''. Two examples: \textit{Creationism} (religious rejection of science) and \textit{Materialist Atheism} (Dawkins: science disproves God).
  \item \textbf{Independence}: ``Strangers'' -- separate compartments, different questions. Science: factual ``is''; Religion: moral ``ought'' / ultimate meaning. \textbf{NOMA} = best example.
  \item \textbf{Dialogue}: ``Limited Partners'' -- discuss boundary questions and methodological similarities. No deep substantive merger.
  \item \textbf{Integration}: ``Substantive Partners'' -- unified metaphysics. Three types: (a) Natural Theology, (b) Theology of Nature, (c) Systematic Synthesis.
\end{enumerate}

\textbf{ECKLUND} (sociologist, 9,400--22,500 scientists globally):
\begin{itemize}[noitemsep]
  \item Independence = most popular worldwide (Taiwan: 62\%)
  \item Western (UK/USA/France): Independence $>$ Conflict $>$ Collaboration
  \item Non-Western (HK, Taiwan, India): Independence $>$ Collaboration (science + religion support each other in education)
  \item UK: highest Conflict (35\%). Science may NOT secularize individuals but CAN secularize broader culture.
\end{itemize}

\vspace{3pt}

\subsection{HOOFDSTUK 3 (vervolg): NOMA en Kritiek}

\textbf{GOULD'S NOMA} (Non-Overlapping Magisteria):

\begin{tabular}{ll}
\toprule
Science Magisterium & Religion Magisterium \\
\midrule
Empirical realm -- ``what is?'' & Ultimate meaning, moral value -- ``what ought?'' \\
Factual character of nature & Human purpose, values \\
\bottomrule
\end{tabular}

Requires \textit{Mutual Humility}: religion refrains from factual claims; science avoids dictating meaning. \textbf{Probleem:} NOMA itself is a normative claim (``should'') $\to$ belongs to the Magisterium of Value, not Science.

\textbf{BURMS \& DE DIJN} -- 3 irreducible human interests:
\begin{enumerate}[noitemsep]
  \item \textit{Manipulative interest}: striving to change reality (technology)
  \item \textit{Cognitive interest}: striving to understand reality (science)
  \item \textit{Search for Meaning}: striving for connectedness, spirituality (religion)
\end{enumerate}
``The heart does not always follow the head'' -- insight $\ne$ experience. A brain scan of someone in love cannot capture the subjective essence of love. $\to$ Meaning is irreducible to cognition.

\textbf{SAM HARRIS} (Scientism): Science CAN determine values because well-being is a mental state $\to$ domain of natural science. Differs from Gould: Harris dissolves NOMA by bringing ethics into science.

\textbf{PETER ATKINS}: All ``why-questions'' are either disguised ``how-questions'' or meaningless.

\textbf{JERRY COYNE} (\textit{Faith vs.\ Fact}, Conflict model): Science and religion are incompatible due to \textit{conflicting methods}. Religion uses faith/authority; science uses reason/evidence. \textbf{6 characteristics of science:} Falsifiability, Doubt/Critical Attitude, Replication, Parsimony (simplest explanation), Living with Uncertainty, Collectivity. \textbf{3 characteristics of religion:} (1) Theism, (2) Moral System (God approves/disapproves), (3) Personal Relationship with God. Coyne's critique of NOMA: ``\textit{homeopathic dilution of religion}'' -- NOMA reduces religion to a version so weakened it no longer resembles real faith.

\vspace{3pt}

\subsection{HOOFDSTUK 3 (\texorpdfstring{p.$\sim$130}{p.~130}): Ichneumon Wasps}

\textbf{Schrijf op/bij de sectie over de ichneumon wasps:}

\textbf{ICHNEUMON WASPS} (parasitoid insects): lay eggs inside living caterpillars; larvae eat host from the inside while keeping it alive. Darwin: ``I cannot persuade myself that a beneficent and omnipotent God would have designedly created the Ichneumonidae with the express intention of their feeding within the living bodies of caterpillars.'' Important for science-religion debate: demonstrates \textit{amoral nature} -- nature is neither good nor bad, undermining the image of a benevolent Creator. Supports the Conflict model (God vs.\ evil in nature). \textbf{Studforum:} appears on exam Sept 2022!

\vspace{3pt}

\subsection{HOOFDSTUK 4 (\texorpdfstring{p.$\sim$192}{p.~192}): Metafysische Ambiguïteit}

\textbf{METAFYSISCHE AMBIGUÏTEIT van moderne wetenschapsbevindingen} (p.192):

The findings of modern natural science are metaphysically ambiguous: the same scientific facts are open to different and conflicting \textit{metaphysical interpretations} (theistic vs.\ atheistic) that cannot be settled by the facts themselves. People with conflicting metaphysical presuppositions look at the same facts through ``glasses of a different colour.'' Illustrated by the \textbf{Parable of the Invisible Gardener}: two explorers find a garden -- one sees evidence of a gardener (God), the other does not. Both share the same empirical data but draw opposite metaphysical conclusions. Also linked to \textbf{Sardar's} more radical view: modern natural science itself approaches reality in a particular (Western) way, so the very framework of science is not metaphysically neutral.

\vspace{3pt}

\subsection{HOOFDSTUK 4 (vervolg): Gaia Hypothese \& Doolittle}

\textbf{GAIA HYPOTHESE} (James Lovelock): Earth is a self-regulating system that maintains conditions for life (as if it were a living organism). Metaphor: Earth as one large organism. Spirituele relevantie: Gaia as an organicist worldview (universe as living whole), contrast with the mechanistic worldview.

\textbf{DAISYWORLD}: Simulation with black and white daisies that regulate Earth's temperature via feedback mechanisms, even as solar intensity increases $\to$ proves Gaia-like regulation can emerge through selfish/evolutionary processes, without teleology. Eventually the sun becomes too intense and the daisies die.

\textbf{FORD DOOLITTLE's KRITIEK op Lovelock}: Natural selection works at the level of individual organisms or genes, not at the level of the whole planet. There is no mechanism by which ``the good of the planet'' can be selected for. Gaia would require organisms to ``altruistically'' sacrifice fitness for planetary regulation -- which evolution cannot produce.

\textbf{REBUTTAL (verdediging van Gaia, aug 2022 examenvraag):} Lovelock's later Daisyworld model shows Gaia-like regulation can emerge from \textit{selfish} individual behaviour, requiring no planetary altruism. The regulation is an \textit{emergent property} of individual interactions. One can also distinguish between Gaia as a \textit{scientific hypothesis} (testable: does Earth regulate conditions?) versus Gaia as a \textit{religious/philosophical metaphor} (Gaia as source of spirituality and meaning).

\vspace{3pt}

\subsection{HOOFDSTUK 5 (§5.3, \texorpdfstring{p.$\sim$165--180}{p.~165--180}): Welcome to the Anthropocene / Spirituele Uitdaging}

\textbf{Schrijf bij begin hfst 5:} $\to$ \textbf{SLEUTELVRAAG 1: ecologische crisis = spirituele crisis}

The Anthropocene represents a \textit{spiritual challenge} (not merely technical): the question of what it means to be human, what our responsibility is, what stories we must tell. The ecological crisis forces us to rethink our place in the cosmos.

\textbf{ECOMODERNISME} (technologisch optimisme): Humans can use technology to ``decouple'' from nature and solve the crisis. Critique in book: this misunderstands the depth of the crisis as a crisis of values.

\textbf{DEEP ECOLOGY}: Humans have no special status; all living beings have equal intrinsic value. Critique: devalues human responsibility.

\vspace{3pt}

\subsection{HOOFDSTUK 6: Vergroening van de Wereldreligies}

\textbf{Schrijf op de eerste pagina van hfst 6 (SLEUTELVRAAG 3 + korte definitie topics):}

\textbf{DRIE MODELLEN IN WESTERSE RELIGIES (theïstisch):}
\begin{enumerate}[noitemsep]
  \item \textbf{Dominion} (heerschappij): humans rule over nature (rejected by most ecotheologians)
  \item \textbf{Stewardship} (rentmeesterschap): humans care for God's creation $\to$ CENTRAL IDEA of greening Abrahamic religions
  \item \textbf{Creation as Community}: humans belong to the community of creation, not above it
\end{enumerate}

\textbf{JODENDOM (p.$\sim$245--255):}

\textbf{Tikkun Olam} (repairing/completing the world): Jewish concept from the Mishnah (protecting the vulnerable), developed by 16th-c. mystic Isaac Luria into cosmic repair: divine vessels shattered at creation, leaving \textit{kelipot} (broken shards trapping light); humans must perform good deeds to release sparks $\to$ ecological application: repairing the world includes ecological repair. Also: daily \textit{Alenu} prayer ends with hope to ``restore creation under God's nurturing role.''

\textbf{Bal Taschit} (do not destroy): prohibition against destroying or wasting anything potentially useful. Rooted in Deuteronomy 20: do not cut down fruit trees during a siege (\textit{Pshat} = literal meaning). \textit{Drash} (rabbinical inquiry) expands this into a universal ban on waste. Ecological principle: stewardship requires non-destruction.

\textbf{Humans as \textit{Adamah} (earth/ground)}: Rabbi Lawrence Troster -- humans are ``earthlings'' made of earth, sharing DNA with other creatures, ``absolute latecomers'' in evolution $\to$ curbs human arrogance.

\textbf{ISLAM (p.$\sim$255--262):}

\textbf{Khalifa} (regent/steward): In Islam, humans are placed on Earth as vicegerents/stewards (\textit{khalifa}) to whom God's property is entrusted. A Khalifa is accountable to God through \textit{Akhrah} (accountability/judgement day). Stewardship rooted in \textbf{Tawhid} (absolute unity of God reflected in the unity of creation): to harm the environment is to violate God's singular design. Assisi Declarations (1986): Abdullah Omar Nasseef -- humans are ``between the angels and the animals'' with the ``ability to think and rebel.''

\textbf{Green Deen/Din} (Tatiana Jerbouh guest lecture): 6 principles of Islamic eco-spirituality. Principles also in book: \textit{Tawhid} (unity), \textit{Khalifa} (stewardship), \textit{Akhrah} (accountability), \textit{Mizan} (balance), \textit{Fitra} (natural disposition), \textit{Halal} (permissible, also in ecology).

\textbf{CHRISTENDOM -- LAUDATO SI' (p.254--262 + 320--327):}

\textbf{Laudato Si'} (2015 Papal Encyclical, Pope Francis): ``Care for Our Common Home.'' Advocates \textbf{Moderate Anthropocentrism} [humans have unique value AND responsibility to care for the ecosystem], rejecting both (1) ``tyrannical'' dominion/exploitation and (2) radical \textbf{Biocentrism} [all living beings have equal inherent value -- feared to devalue human responsibility]. Draws on \textbf{St.\ Francis of Assisi} (patron saint of ecologists): ``Retreat into Creation'' -- being with other creatures as brothers and sisters. \textbf{Moderate Anthropocentrism is based on THEOCENTRISM}: God is the centre $\to$ humans are unique (created in God's image) but as stewards, not lords; value comes from God, not from human dominance. Tragedy: political polarisation means $\sim$42\% of American Catholics allow political preferences to override this teaching.

\textbf{OOSTERSE RELIGIES (p.$\sim$262--280):}

\textbf{HINDOEÏSME}: \textit{Sarvatma-bhava} (everything is connected), \textit{Svarupa} (concrete embodiment of the divine, e.g.\ a sacred river), \textit{Seva} (loving service/environmental activism), \textit{Sambandha} (firm relationship with divine-in-nature). Bhakti Yoga (path of loving devotion).

\textbf{BOEDDHISME}: \textbf{Paticca-samuppada} (dependent origination / radical interconnectedness, p.265--266): everything arises in dependence on conditions; there is no permanent separate self. Eco-Buddhists (Joanna Macy): aligns with cybernetics/systems thinking. Transition from \textit{Ego-self} to \textit{Eco-self} (self co-extensive with all beings).

\textbf{CONFUCIANISME}: \textbf{Ren} (p.267--268): humaneness / benevolence -- ``being one body with all things''; the heart-mind of humans serves as the centre of the cosmic body. Humans are not separate from the cosmos but its heart.

\textbf{TAOÏSME}: \textbf{Wuwei} (p.269): non-action, non-interference with natural processes. Expressed in \textit{Shanshui hua} (landscape paintings of ``mountains and water''): accepting and flowing with natural change rather than dominating it.

\textbf{Multiple Religious Affiliation} (Japan, East Asia): participating in rituals of several religions simultaneously. \textit{Orthopraxis} (right practice/ritual) prioritised over exclusive Western ``belief.''

\vspace{3pt}

\subsection{HOOFDSTUK 7 (§7.1, §7.2.3, p. ca. 273--310): Aardse Systeemwetenschap als Spirituele Bron}

\textbf{Schrijf bij begin hfst 7:} $\to$ \textbf{SLEUTELVRAAG 2b: Spiritualiteit voor het Anthropoceen}

\textbf{EARTHRISE} (p.273--274): Photograph taken by Apollo 8 crew (December 1968) showing Earth rising above the moon's surface. First time humanity saw Earth from outside. Effect: awareness of Earth as fragile, united, alone in space. Source of spirituality: connectedness (we are all on this tiny blue dot), self-transcendence (we are small), meaning (we must care for this home). Connected to \textbf{The Blue Marble} (1972, Apollo 17) and \textbf{Pale Blue Dot} (Carl Sagan, Voyager 1, 1990).

\textbf{AXIAL/POST-AXIAL} (from guest lecture, June 2025 exam): \textit{Axial Age} (Karl Jaspers): 800--200 BCE, simultaneous emergence of reflective, transcendent religious/philosophical traditions across cultures (Buddha, Confucius, Greek philosophers, Hebrew prophets). Post-axial: religions that developed after, building on axial traditions.

\textbf{§7.2.3 Spiritualiteit en Natuurwetenschap:}
Three spiritual functions of science: (1) Awe and wonder at complexity of nature $\to$ connectedness; (2) Humility: science reveals we are not the centre (``Crowning Glory'' critique: humans are not pinnacle but a ``tiny twig'' on a bush dominated by bacteria) $\to$ self-transcendence; (3) The universe's story as a source of meaning $\to$ narrative/meaning.

\vspace{3pt}

\subsection{HOOFDSTUK 8 (§8.2): Environmentalisme als Religie}

\textbf{Schrijf bij begin hfst 8:} $\to$ \textbf{SLEUTELVRAAG 4: Is environmentalisme een religie?}

\textbf{MICHAEL CRICHTON's stelling:} ``Environmentalism is the religion of choice for urban atheists.'' The structure of environmentalism mirrors religion: (1) Eden = pristine nature; (2) Fall = pollution, industrialisation; (3) Judgment Day = climate catastrophe; (4) Salvation = sustainability. ``Sustainability is salvation in the Church of the Environment.'' Problem: facts are irrelevant because beliefs are about faith, not evidence.

\textbf{JOEL GARREAU:} Environmentalism as a religion -- characterised by belief, salvation narrative, moral community, and a sense of sin (carbon sins). Considers this a \textit{problem}: without empirical grounding, environmentalism becomes dogmatic.

\textbf{RELIGIEDEFINITIES -- toetsen:} Depending on how you define religion: (a) if religion requires theism/ritual/community $\to$ environmentalism fails; (b) if religion is a meaning-giving worldview with moral norms and a narrative of Fall/Salvation $\to$ environmentalism qualifies. Book's nuanced answer: environmentalism shares \textit{functional} features of religion (provides meaning, community, moral framework, sacred/profane distinctions) but lacks key \textit{substantive} features (explicit theism, transcendent reference).

\textbf{LINK MET BOEK:} Religious vs.\ irreligious in environmentalism -- some religious people reject environmentalism (White critique), others embrace it (Laudato Si', Tikkun Olam). Secular environmentalism can have quasi-religious features.

\vspace{8pt}

%=========================================================
\section{VOLLEDIGE ANTWOORDEN OP DE SLEUTELVRAGEN}
%=========================================================

\textbf{Schrijf deze antwoorden IN HET BOEK bij de relevante hoofdstukken (zie pijlen hierboven). Ze zijn max 2 pagina's lang op het examen.}

\vspace{3pt}

\begin{answerbox}
\subsection{SLEUTELVRAAG 1: Is de ecologische crisis een spirituele crisis? (Hfst 1--2)}

\textit{[Schrijf bij hfst 1--2 beginpagina]}

The ecological crisis is fundamentally a \textbf{spiritual crisis} -- a crisis of meaning, values, and our self-understanding as humans. The book argues that material factors (overpopulation, pollution, technology) are merely \textit{symptoms}; the root cause is a profound shift in how we understand our place in the cosmos.

\textbf{Evidence (from book):}
(1) Lynn White's thesis: the crisis has \textit{religious roots} in Western Christianity's desacralization of nature. By removing spiritual significance from nature and declaring humans its lords, Western Christianity provided the ideological basis for exploitation. This shows the crisis is not merely technical but stems from a worldview.

(2) The crisis raises irreducibly spiritual questions: What is the human being? What is our responsibility to the Earth? What stories must we tell? These are questions of \textit{meaning} (Burms and De Dijn's third irreducible interest), not just of cognition or manipulation.

(3) Ecomodernism (the view that technology can solve the crisis through decoupling) fails because it ignores the spiritual dimension: even if we could technically fix the climate, we would not have addressed the underlying cultural attitude of domination and exploitation.

\textbf{Assessment:} I agree with the book's diagnosis. Engineering education often focuses on technical solutions, but the root problem is an unsustainable way of relating to the world. Laudato Si' (Francis, 2015) also diagnoses a \textit{technocratic paradigm} -- a worldview in which everything is treated as a technical problem to be solved, mastered and exploited. A spiritual transformation -- a new understanding of our place in the cosmos -- is needed alongside technical innovation. Guest lectures (Tiqqun Olam, Green Islam, Bah\'a'\'i\xspace moderation) all confirm: the solution must include a transformation of values.

\textbf{Counter:} One could argue the crisis is primarily material (greenhouse gas concentrations, deforestation rates) and can be solved materially. But: even if true, the motivation to solve the crisis requires a moral and spiritual commitment that material science cannot provide. Therefore, spirituality remains necessary.
\end{answerbox}

\vspace{3pt}

\begin{answerbox}
\subsection{SLEUTELVRAAG 2a: Relatie Wetenschap en Religie (Hfst 3--4)}

\textit{[Schrijf bij hfst 3, p.76--80]}

The book elaborates several models of how science and religion relate, built on \textbf{Ian Barbour's} typology.

\textbf{Conflict model} (Coyne, Dawkins): science and religion make incompatible claims. Religion relies on faith/authority; science on evidence/falsifiability. The Ichneumon wasp case: amoral nature contradicts a benevolent Creator God. Assessment: the conflict is real for \textit{ordinary} religion (which does make factual claims), but the model is too simplistic -- it ignores the possibility of different domains.

\textbf{Independence model -- NOMA} (Gould): science covers empirical facts (``is''), religion covers ultimate meaning and moral values (``ought''). This is the most popular model globally (Ecklund). Assessment: philosophically attractive, but (1) Coyne's critique: NOMA requires a ``homeopathic dilution'' of religion -- ordinary believers DO make factual claims that science contradicts; (2) The NOMA principle itself is normative, not scientific. (3) Metaphysical ambiguity (p.192): same facts admit multiple interpretations, but science cannot adjudicate between them.

\textbf{Dialogue/Integration} (Barbour): science and religion genuinely enrich each other. Natural Theology, Theology of Nature, Systematic Synthesis. Gaia hypothesis (Lovelock) as example of integration: science and an organicist spiritual vision of Earth as living whole.

\textbf{My view}: I lean toward \textbf{Independence with dialogue}. Science and religion operate in different domains (fact vs.\ meaning), but they cannot be completely sealed off: Burms and De Dijn show that the search for meaning is an irreducible human interest that science cannot satisfy, yet must coexist with scientific knowledge. The Bah\'a'\'i\xspace view is illuminating: ``Religion without science becomes superstition; science without religion becomes crude materialism.''
\end{answerbox}

\vspace{3pt}

\begin{answerbox}
\subsection{SLEUTELVRAAG 2b: Spiritualiteit voor het Anthropoceen (Hfst 7)}

\textit{[Schrijf bij hfst 7, §7.1--7.2.3]}

A spirituality for the Anthropocene must address our situation as a species that has become a geological superpower, yet remains morally and spiritually immature.

\textbf{Key elements (from book):}

(1) \textbf{Earth system science as spiritual source} (§7.2.3): The Earthrise photograph (1968) and Pale Blue Dot (Sagan) generated genuine spiritual experiences: awe, connectedness, humility. These are not religious in the traditional sense but fulfil the three attributes of spirituality (connectedness, self-transcendence, meaning). A spirituality for the Anthropocene can build on this.

(2) \textbf{Gaia} (Lovelock): Earth as a self-regulating system -- not literally alive, but functioning as a whole. An organicist vision (cosmos as living whole) rather than mechanistic (cosmos as machine). This generates a sense of belonging to something larger. The Bah\'a'\'i virtue of \textit{moderation} also fits: recognising our dependence on Earth ``beneath our feet'' should cleanse us of pride and arrogance.

(3) \textbf{Eco-self} (Buddhism, Macy): the transition from ego-self to a self co-extensive with all beings. Paticca-samuppada: radical interconnectedness. Practically: caring for Earth = caring for the self.

(4) \textbf{Multiple religious resources}: Tikkun Olam (repair the world), Khalifa (stewardship and accountability), Laudato Si' (moderate anthropocentrism, fraternal love for creation), Ren (cosmic body), Wuwei (non-interference, flow with nature).

\textbf{My view}: A spirituality for the Anthropocene must be \textit{open-source}: drawing on scientific awe, religious traditions of stewardship and interconnectedness, and philosophical humility about our place. It must avoid two errors: (a) \textit{triumphalism} -- thinking technology will save us (ecomodernism); (b) \textit{nihilism} -- paralysis in the face of the crisis. As an engineer, the insight that work can be a form of service (Bah\'a'\'i) and that every design choice participates in a broader vision of humanity is personally meaningful.
\end{answerbox}

\vspace{3pt}

\begin{answerbox}
\subsection{SLEUTELVRAAG 3: Rol van Wereldreligies bij Ecologische Crisis (Hfst 6)}

\textit{[Schrijf bij hfst 6, begin]}

\textbf{Potential roles -- a nuanced position:}

\textbf{Arguments for a significant role:}
(1) Religions provide \textit{moral motivation}: Tikkun Olam (Judaism), Khalifa (Islam), Laudato Si' (Catholicism), and Paticca-samuppada (Buddhism) all generate strong ecological obligations rooted in deep spiritual commitments -- more durable than mere policy preferences.
(2) Scale: religions involve billions of people. If Laudato Si' mobilised even 10\% of 1.3 billion Catholics, this would be a massive force.
(3) Narrative and meaning: religions provide the stories and rituals that make ecological commitment emotionally sustainable.
(4) Guest lectures confirm: Tiqqun Olam (repairing the world) and Green Islam (Khalifa, Tawhid, Akhrah) are being actively applied to ecological activism.

\textbf{Arguments for limitations:}
(1) White's thesis: religions (Western Christianity especially) helped \textit{cause} the crisis. Internal reform takes time.
(2) Political polarisation: Laudato Si' shows that $\sim$42\% of American Catholics ignore it due to political identity.
(3) Institutional conservatism: religions can resist change; their official positions may lag behind ecological urgency.
(4) Exclusivism: some religious frameworks require belief in God -- those who are non-religious need secular ecological motivation.

\textbf{Synthesis (my position):} World religions CAN play a significant positive role, but only if they (1) engage in genuine self-critique (confession of guilt, not just apologetics), (2) move beyond internal theological debates to concrete ecological action, and (3) cooperate across traditions (as in the Assisi Declarations). Crucially, the insights of religion can also be ``secularised'' -- the concept of stewardship, or of radical interconnectedness, can motivate environmentalists regardless of whether they are religious. Religion is not the only path, but it remains a uniquely powerful one.
\end{answerbox}

\vspace{3pt}

\begin{answerbox}
\subsection{SLEUTELVRAAG 4: Is Environmentalisme een Religie? (Hfst 8)}

\textit{[Schrijf bij hfst 8, §8.2]}

\textbf{Arguments that it IS a religion (Crichton, Garreau):}
\begin{itemize}[noitemsep]
  \item Narrative structure: Eden (pristine nature) $\to$ Fall (pollution) $\to$ Judgment Day (climate catastrophe) $\to$ Salvation (sustainability)
  \item Sacred/profane distinction: ``carbon sinners'' vs.\ the saved
  \item Moral community with shared beliefs not dependent on empirical scrutiny
  \item Rituals (recycling, Fridays for Future) and prophets (Greta Thunberg)
  \item Provides meaning, identity, and community $\to$ functional definition of religion satisfied
\end{itemize}

\textbf{Arguments that it is NOT a religion:}
\begin{itemize}[noitemsep]
  \item Lacks theism and a transcendent reference point
  \item Lacks formal ritual, sacred texts, or clergy in the traditional sense
  \item Is empirically grounded: claims about climate change are falsifiable and evidence-based $\to$ meets Coyne's criteria for science, not faith
  \item Depends on scientific consensus, not revelation or authority
\end{itemize}

\textbf{Nuanced conclusion}: Whether environmentalism is a religion depends on the definition of religion. Under a \textit{substantive} definition (requires theism/transcendence), environmentalism is NOT a religion. Under a \textit{functional} definition (provides meaning, moral norms, community, salvation narrative), environmentalism IS religion-\textit{like}. The book's view: environmentalism shares functional features of religion but is not identical. This is important: treating environmentalism as a religion (as Crichton does) can be used to \textit{dismiss} it as irrational dogma, ignoring its empirical foundation. But recognising its quasi-religious function can help us understand why it motivates people so powerfully.
\end{answerbox}

\vspace{3pt}

\begin{answerbox}
\subsection{SLEUTELVRAAG 5: Technologie, Religie en Spiritualiteit (p.343--351, §2.3.1, §2.3.3)}

\textit{[Schrijf bij hfst 5 / p.343--351]}

\textbf{Barbour/Niebuhr's 5 Christ-and-Culture models applied to technology:}
\begin{enumerate}[noitemsep]
  \item Salvation \textit{outside} technology (Christ against culture) -- Ellul: \textit{technique} (rationalised pursuit of efficiency) dominates; technology enslaves; salvation must come from outside. Tech-pessimist.
  \item Salvation \textit{by means of} technology (Christ of culture) -- Dessauer, Hefner: technology as participation in God's creation; humans as ``created co-creators.'' Tech-optimist.
  \item Salvation \textit{beyond} technology (Christ above culture) -- Aquinas: technology is good but insufficient; grace needed. Tech is preparatory.
  \item Salvation and technology \textit{in paradox} (Luther): technology belongs to fallen world; necessary but dangerous.
  \item Salvation \textit{requires conversion} of technology (Augustine, Schuurman): technology must be redirected toward God's purposes.
\end{enumerate}

\textbf{Secular equivalents:}
\begin{itemize}[noitemsep]
  \item \textit{Techno-pessimism} (Ellul, Kaczynski): technology threatens humanity
  \item \textit{Techno-optimism} (Kurzweil): Singularity, merging with AI, defeating death $\to$ secular salvation. David Noble: ``Religion of Technology'' -- technology promises immortality/control/liberation (Apple 1984 ad).
  \item \textit{Pope Francis} (realistic): technology is good (medicine, engineering, sustainable innovation) but becomes dangerous when guided by the \textit{technocratic paradigm} (treating everything as a technical problem; mastery as dominant mode of relating to reality).
\end{itemize}

\textbf{Miss Case (p.65):} Technology as a source of spirituality: by connecting us and giving access to vast information, technology enables new forms of connectedness and self-transcendence -- attributes of spirituality. But: this only works if technology is used reflectively, not addictively.

\textbf{My view:} Technology and spirituality are neither simply opposed nor simply identical. Technology has genuinely spiritual functions (connectedness via the internet, awe at scientific instruments revealing the cosmos) but also spiritual dangers (technocratic reduction of reality to what is measurable and manipulable). The Bah\'a'\'i\xspace principle -- science without ethics becomes crude materialism; religion without science becomes fanaticism -- captures the needed balance. As a future engineer: every design participates in a vision of humanity. This course has taught me that technical choices are always also moral and spiritual choices.
\end{answerbox}

\vspace{8pt}

%=========================================================
\section{VRAAG 3: PERSOONLIJKE REFLECTIE (TEMPLATE)}
%=========================================================

\textit{[Schrijf dit op de achterste pagina's van het boek -- klaar om te kopiëren!]}

\begin{answerbox}
\textbf{REVIEW QUESTION TEMPLATE (schrijf in het boek, pas aan op het examen)}

The most important insight I take from this course is that the ecological crisis is not merely a technical challenge but a spiritual one: a crisis of the stories we tell ourselves about who we are and what we are for. Before this course, I understood the Anthropocene primarily through the lens of engineering solutions -- solar panels, carbon capture, efficiency gains. The course has challenged me to see that even if these solutions worked technically, they would not address the underlying attitude: the \textit{technocratic paradigm} (Pope Francis) that treats nature as raw material to be mastered.

What surprised me most was the concept of \textit{metaphysical ambiguity} (p.192): the same scientific facts -- the Big Bang, evolution, the Earthrise photograph -- admit both theistic and atheistic interpretations. Science cannot resolve the deepest questions of meaning. This explains why Burms and De Dijn argue that the search for meaning is irreducible to cognition: a perfect brain scan of someone in love cannot capture what it means to be in love. I found this genuinely convincing.

I disagreed most strongly with strict NOMA (Gould). While I appreciate the goal of peaceful coexistence, Coyne's ``homeopathic dilution'' critique resonates: a religion that makes no factual claims about the world is unrecognisable to billions of actual believers. In practice, science and religion DO overlap, and the tensions are real and productive, not to be simply partitioned away.

As an engineer, the course has permanently changed how I think about my profession. The Baha'i insight -- that work can be a form of service/worship -- and the Christian notion of ``created co-creators'' (Hefner) both suggest that engineering is never morally neutral. Every bridge, every circuit board, every algorithm participates in a broader vision of what humanity is for.
\end{answerbox}

\vspace{8pt}

%=========================================================
\section{GUEST LECTURES -- SAMENVATTING (schrijf op lege pagina's achteraan)}
%=========================================================

\subsection{Guest Lecture 1: Tiqqun Olam \& Bal Taschit (Joodse ecologie)}
\textbf{Tikkun Olam}: Repairing the world -- cosmic and social repair. Good deeds release trapped divine sparks (\textit{kelipot}). Ecological application: environmental activism as repair. \textbf{Bal Taschit}: Do not destroy. Universal ban on waste. \textbf{Link to book}: matches the \textit{stewardship} model and \textit{creation as community}; matches Adamah (humans as earthlings). Engineers: building sustainably = Tikkun Olam.

\subsection{Guest Lecture 2: Green Islam (Tatiana Jerbouh)}
\textbf{6 principles of Green Deen:} (1) Tawhid (unity of God/creation), (2) Khalifa (stewardship), (3) Akhrah (accountability), (4) Mizan (balance -- do not exceed limits), (5) Fitra (innate ecological nature of humans), (6) Halal (permissible also ecologically). \textbf{Link to book}: Khalifa (p.$\sim$255--260), Tawhid, Akhrah all in ch.6. Same principles as book's Islam section. Assisi 1986: Omar Nasseef.

\subsection{Guest Lecture 3: Bah\'a'\'i Faith (Science, Ecology, Spiritual Development)}
\textbf{Key ideas:} Progressive Revelation (God sends messengers adapted to humanity's maturity). Unity: of God, of Religion, of Humanity. Harmony of science and religion: ``religion without science = superstition; science without religion = crude materialism.'' \textbf{Moderation}: technology and economy must not exceed ecological limits. Work as Worship: technical work is spiritually meaningful when performed as service. Independent Investigation of Truth: faith must be personally examined, not inherited blindly. \textbf{Link to exam key question 5}: technology, religion and spirituality.

\subsection{Guest Lecture 4: Christian Perspectives on Technology (13 maart 2026, Depoortere)}
\textbf{H. Richard Niebuhr -- \textit{Christ and Culture} (1951):} 5 modellen van hoe christenen omgaan met ``de wereld''; hier toegepast op technologie. Cultuur = menselijke inspanning om de kosmos te ordenen (cf. Gen 1:27-28, Cultural Mandate).

\textbf{Niebuhr's 5 modellen $\to$ Salvation \& Technology:}
\begin{enumerate}[noitemsep]
  \item \textbf{Christ against culture} $\to$ \textbf{Salvation Outside Technology} (Ellul, Blair): technologie is kwaad, scheidt je ervan. Seculiere versie: techno-pessimisten (Kaczynski: ``The Industrial Revolution has been a disaster'').
  \item \textbf{Christ of culture} $\to$ \textbf{Salvation By Means of Technology} (Hefner, Dessauer): technologie als deelname aan Gods schepping; ``created co-creators.'' Seculiere versie: techno-optimisten/enthusiasten.
  \item \textbf{Christ above culture} $\to$ \textbf{Salvation Beyond Technology}: technologie is goed maar onvoldoende; genade nodig.
  \item \textbf{Christ and culture in paradox} $\to$ \textbf{Salvation and Technology in Paradox}: technologie behoort tot de gevallen wereld; noodzakelijk maar gevaarlijk.
  \item \textbf{Christ transformer of culture} $\to$ \textbf{Salvation Requires Conversion of Technology} (Schuurman): technologie moet omgebogen worden naar Gods doelen. $\Rightarrow$ \textbf{Technology is in need of conversion.}
\end{enumerate}

\textbf{Menselijke kosten van technologie (Ellul, 9 punten):} (1) uniformiteit in massamaatschappij; (2) enge focus op rationaliteit/efficiëntie; (3) onpersoonlijkheid en manipulatie; (4) risico op technologie die oncontroleerbaar wordt; (5) vervreemding van arbeiders; (6) link technologie--militair; (7) milieuproblemen; (8) verslechtering intermenselijk contact; (9) instorting gemeenschapsleven.

\textbf{Fundamentele spanning:} Kurzweil (``science/tech has replaced religion'' -- The Singularity is Near/Nearer) VS David Noble (``science and technology ARE the religions of the modern age'' -- \textit{The Religion of Technology}). Jeffrey C. Alexander: inventies als stoomkracht, spoor, telegraf werden ``hailed as vehicles for secular transcendence''; ingenieurs verheven tot ``worldly priests.'' Apple 1984 ad = technologie als verlossing?

\textbf{Realism about technology (Pope Francis / Laudato Si' ch.3):} 4 vereisten: (1) niet verwachten dat technologie de crisis in een oogwenk oplost; (2) niet bezwijken voor escapistische fantasieën; (3) niet blind zijn voor hoe de crisis het huidige systeem (Science-Technology-Culture) fundamenteel in vraag stelt; (4) beseffen dat technologische innovatie veel tijd kost -- tijd die we misschien niet meer hebben. LS §102--114: Technology: Creativity \& Power + Globalization of the Technocratic Paradigm.

\textbf{Egbert Schuurman quote:} ``Let me conclude by presenting technology as a pilgrimage of obedience, a mandated way to greater insight into the meaning of creation as the Kingdom of God. We are called upon to honor the Lord in technology as in every sphere of life.''\textbf{ $\to$ technologische bekering als antwoord op de spanning.}

\subsection{Guest Lecture 5: Existential Challenges of Planetary Health (Tom Hannes)}
\begin{exambox}
\textbf{EXAMEN JUNI 2025 -- DIRECT GEVRAAGD (Q1 a/b/c):}
\begin{itemize}[noitemsep]
  \item \textbf{a. What is Blue Marble?} Foto genomen door Apollo 17 crew (december 1972) op weg naar de maan. Eerste foto waarop de gehele verlichte Aarde te zien is. Anders dan Earthrise (Aarde \textit{opkomend boven de maan}) toont Blue Marble de Aarde als een volledig, prachtig blauw geheel in de zwarte ruimte.
  \item \textbf{b. Kritiek van Tom Hannes op Blue Marble:} De Blue Marble foto genereert een ``overview effect'' -- een ervaring van de Aarde als een uniek, fragiel, prachtig geheel. Maar dit kan leiden tot een \textit{abstract, detached} beeld van de planeet: men ziet de Aarde van buiten als een object, niet van binnenuit als een levende gemeenschap waarvan men deel uitmaakt. De afstandelijkheid van de ``view from above'' staat op gespannen voet met de \textit{incarnate}, lokale, belichaamde verbondenheid met de natuur die ecologisch engagement vereist. Ook: de foto kan een gevoel van eenheid suggereren dat politieke en culturele spanningen maskeert.
  \item \textbf{c. Axial vs.\ Post-axial (Tom Hannes):} \textbf{Axiale periode} (Karl Jaspers, 800--200 v.Chr.): een tijdperk van gelijktijdige, onderling onafhankelijke religieuze en filosofische doorbraken wereldwijd -- Boeddha, Confucius, Griekse filosofen (Socrates/Plato), Hebreeuwse profeten -- gekenmerkt door: (a) zelfreflectie en kritisch denken, (b) transcendentie (verwijzing naar een werkelijkheid buiten het gewone), (c) universalisme. \textbf{Post-axiale religies}: religies die na de Axiale Periode zijn ontstaan en voortbouwen op axiale tradities (bv.\ Christendom, Islam bouwen voort op Joods-axiaal denken). Tom Hannes: de Earthrise- en Blue Marble-foto's kunnen een \textit{nieuwe axiale ervaring} inluiden -- een transformatie van bewustzijn voor de Antropoceen-generatie die de Aarde als fragiele, kwetsbare gemeenschap ziet.
\end{itemize}
\end{exambox}

\textbf{Earthrise (1968) vs.\ Blue Marble (1972):}
\begin{itemize}[noitemsep]
  \item \textbf{Earthrise} (Apollo 8, december 1968): Aarde \textit{opkomend} boven de maanhorizon. Eerste blik op Aarde van buiten. Gave een gevoel van eenheid, kwetsbaarheid en wonder. Inspireerde ecologische beweging.
  \item \textbf{Blue Marble} (Apollo 17, december 1972): Volledige, verlichte Aarde als blauw marmer. Meest gereproduceerde foto ooit. Genereert gevoel van Aarde als één kostbaar geheel.
  \item \textbf{Pale Blue Dot} (Voyager 1, Carl Sagan, 1990): Aarde als een minuscuul stofje in een reusachtig heelal. Bron van spirituele awe en bescheidenheid.
\end{itemize}

\textbf{Axiale Periode als spirituele bron:} De Axiale Periode liet zien dat diepe spirituele en filosofische transformaties mogelijk zijn -- ook vandaag is een nieuwe ``axiale'' transformatie nodig om de Anthropoceen-crisis het hoofd te bieden. Tom Hannes: de confrontatie met de planetaire crisis is een uitnodiging tot een nieuwe zelfreflectie en transcendente oriëntatie -- een spiritualiteit die de Aarde herkent als onze enige en gezamenlijke thuis.

\textbf{Link met boek (hfst 7, §7.1--7.2.3, p.$\sim$273--310):} Earthrise als bron van spiritualiteit, Gaia als organistisch wereldbeeld, de spirituele functies van wetenschap (awe, humility, meaning). Blue Marble als \textit{uitbreiding}: de Aarde als object van zorg en verantwoordelijkheid.

\textbf{Planetaire grenzen (in context):} De Aarde heeft fysieke grenzen (Rockström et al.\ planetaire grenzensysteem) -- overschrijding bedreigt de condities voor menselijk en ander leven. Dit is de \textit{existentiële} uitdaging: hoe leven we zinvol in de wetenschap dat we deze grenzen overschrijden?

\subsection{Guest Lecture 6: The Religious and Irreligious in Environmentalism (20 maart 2026)}
\textbf{Kernvraag:} Hoe verhoudt het traditionele milieuactivisme zich tot religie -- en is dit een probleem?

\textbf{Is environmentalisme een religie?} (Link met hfst 8, SV4)
Traditioneel environmentalisme behandelt ``natuur'' als een onveranderlijke, heilige orde waaraan mensen zich moeten aanpassen. Ecomodernisten (Lynas, Brand) noemen dit een \textit{fundamenteel religieuze houding}: ``the traditional green longs to live in harmony with nature'' -- maar dit is, aldus ecomodernisten, een religieus ideaal, geen wetenschappelijk inzicht.
\begin{itemize}[noitemsep]
  \item \textbf{Maarten Boudry:} Ideeën over ``natural order, harmony and Mother Earth'' in de traditionele groene beweging zijn \textit{``harmful nonsense''} $\to$ we moeten religieus denken over natuur vervangen door rationeel denken.
  \item \textbf{Bevestiging vanuit Laudato Si':} LS 68: ``respect the laws of nature and delicate equilibria''; LS 71: ``recovering and respecting the rhythms inscribed in nature by the Creator''; LS 115: ``respect the natural and moral structure [from God]''; LS 117: ``the message contained in the structures of nature itself'' -- dit bevestigt dat zelfs de paus natuur als normatief kader gebruikt.
\end{itemize}

\textbf{Ecomodernistische kritiek op ecologische bekering (ecological conversion):}
\begin{itemize}[noitemsep]
  \item Niemand is bereid zijn levensstandaard op te offeren voor het milieu
  \item Enorme kloof tussen de schaal van het probleem en de schijnbare futiliteit van individuele acties
  \item Zelfbeperking koopt hooguit tijd maar lost niets structureel op
  \item Alle verantwoordelijkheid bij het individu leggen is ineffectief
\end{itemize}

\textbf{Ecomodernistisch alternatief (5 claims):}
\begin{enumerate}[noitemsep]
  \item Er was nooit een paradijs (``There never was a paradise'') -- het verleden was ook geen ecologisch Eden
  \item Moraliseren en schuldgevoel opwekken zijn contraproductief
  \item We need to tell a \textit{positive} story over de toekomst
  \item We hebben \textit{meer} vooruitgang en modernisering nodig, niet minder
  \item We hebben \textit{meer} technologie en innovatie nodig, niet minder
\end{enumerate}

\textbf{Ecomodernisme claimt de Verlichting:} Rationeel denken, geloof in wetenschap, optimisme en hoop, humanisme, materiële groei voor iedereen, liefde voor experiment, verwerping van elke orthodoxie of dogma.

\textbf{Ecomodernistisch Manifest -- Decoupling:} Mensheid moet haar impact op het milieu terugschroeven om meer ruimte voor natuur te maken, MAAR moet niet ``harmoniseren'' met natuur (= religieuze houding). ``Natural systems will not be protected by the expansion of humankind's dependence upon them.'' Decoupling = menselijke materiële afhankelijkheid van natuur verminderen via technologie.

\textbf{Maar decoupling $\ne$ puur rationeel/utilitair leven:} Het Ecomodernistisch Manifest erkent: mensen schrijven dit ``out of deep love and emotional connection to the natural world'' -- spirituele en esthetische waarde van natuur blijft. Decoupling biedt \textit{de mogelijkheid} dat onze materiële afhankelijkheid minder destructief wordt, niet dat we natuur opgeven.

\textbf{Fundamentele spanning (leading back to White):} Is de ecologische crisis \textit{first and foremost} een kwestie van verkeerde mentaliteit (White $\to$ ecologische bekering nodig) of een kwestie van verkeerde technologie (ecomodernisten $\to$ technologische innovatie nodig)? \textbf{Concerns over ecomodernisme:} Geen correlatie tussen religieuze affiliatie en milieu-engagement; geen verschil tussen westerse en niet-westerse contexten in de praktijk; vraag: hebben ideeën echt zoveel impact op gedrag? (cf.\ Vermeersch \& Harari). Maar ecomodernisten blijven geloven dat ideeën tellen.

\textbf{Ecomodernisme en spiritualiteit voor het Anthropoceen:}
\begin{itemize}[noitemsep]
  \item Stewart Brand: \textit{``We are as gods and might as well get good at it.''}
  \item Mark Lynas: \textit{The God Species: Saving the Planet in the Age of Humans}
  \item Yuval Noah Harari: \textit{Homo Deus: A Brief History of Tomorrow}
  \item Vraag: wat betekent ecomodernisme voor onze zoektocht naar spiritualiteit voor het Anthropoceen? Kan een niet-religieuze, technologisch optimistische houding spiritueel bevredigend zijn?
\end{itemize}

\textbf{Link met SV4 (is environmentalisme een religie?) en SV1 (ecologische crisis als spirituele crisis):} Het debat tussen traditioneel environmentalisme (ecologische bekering, stewardship, harmonie met natuur) en ecomodernisme (rationele decoupling, technologische innovatie) loopt parallel aan het debat over of de ecologische crisis primair een spirituele/culturele crisis (White) is of een technische uitdaging. Beide perspectieven zijn nodig: spirituele transformatie van waarden \textit{en} technologische innovatie.

\vspace{8pt}

%=========================================================
\section{DEFINITIE-FLASHCARDS (voor Vraag 1)}
%=========================================================

\textbf{Schrijf deze definities verspreid door het boek bij de betreffende secties:}

\begin{multicols}{2}
\textbf{Tiqqun Olam}: Repairing/completing the world. Jewish concept of cosmic and social repair through good deeds; ecological application: environmental activism. (p.$\sim$245--255, hfst 6)

\textbf{Bal Taschit}: Do not destroy/waste. Jewish ethical principle from Deut.20; universal ban on destroying anything useful. (hfst 6)

\textbf{Khalifa}: Islamic concept of humans as God's stewards/vicegerents on Earth; accountable to God (Akhrah) for treatment of creation. Rooted in Tawhid. (p.$\sim$255--262, hfst 6)

\textbf{Laudato Si'}: 2015 Papal Encyclical by Pope Francis on care for our common home; advocates Moderate Anthropocentrism based on Theocentrism; rejects dominion AND biocentrism. (p.254--262, 320--327)

\textbf{Moderate Anthropocentrism}: Humans have unique value AND responsibility to care for the ecosystem; based on Theocentrism (God's ownership); rejects exploitation AND radical biocentrism. (p.320--327)

\textbf{NOMA}: Non-Overlapping Magisteria (Gould): science covers empirical facts; religion covers ultimate meaning and moral value. Principled and respectful separation. (hfst 3)

\textbf{Paticca-samuppada}: Buddhist dependent origination -- everything arises in dependence on conditions; no permanent separate self; basis for ecological interconnectedness. (p.265--266)

\textbf{Ren}: Confucian humaneness -- being ``one body with all things''; heart-mind of humans as centre of the cosmic body. (p.267--268)

\textbf{Wuwei}: Taoist non-action; non-interference with natural processes; accepting and flowing with natural change. (p.269)

\textbf{Earthrise}: Photograph by Apollo 8 crew (1968); first image of Earth rising above moon's horizon; spiritual source of awe, connectedness, humility; inspired ecological movement. (p.273--274)

\textbf{Blue Marble}: Photograph by Apollo 17 crew (1972); first image of fully illuminated Earth; shows Earth as a complete, fragile blue whole. \textbf{Tom Hannes critique}: generates abstract ``view from above'' that can detach us from embodied, local connection to nature. (GL5 -- \textbf{EXAM juni 2025 Q1a+b!})

\textbf{Axial Age / Post-axial} (Tom Hannes, Karl Jaspers): Axial Period 800--200 BCE: simultaneous emergence of reflective, transcendent traditions (Buddha, Confucius, Greek philosophers, Hebrew prophets): self-reflection, transcendence, universalism. Post-axial: religions building on axial traditions (Christianity, Islam). Tom Hannes: Earthrise/Blue Marble may herald a \textit{new axial experience} for the Anthropocene. (\textbf{EXAM juni 2025 Q1c!})

\textbf{Stewardship}: The central idea of the greening of the Abrahamic religions; humans as caretakers (not owners) of God's creation; reinterprets Genesis ``dominion'' as mandate to maintain balance.

\textbf{Metaphysical ambiguity}: The same scientific findings (e.g.\ evolution, Big Bang) are open to both theistic and atheistic interpretations that cannot be settled by the empirical facts alone. (p.192)

\textbf{Tawhid}: Absolute unity of God reflected in the unity of creation (Islam); harming the environment violates God's singular design.

\textbf{Ecotheologie}: A form of constructive theology focusing on the interrelationships between religion and nature; sparked by White's 1967 thesis.

\textbf{Desacralization}: Removal of spiritual significance from the natural world; White argues Western Christianity desacralized nature.

\textbf{Gaia hypothesis} (Lovelock): Earth as a self-regulating system that maintains conditions for life; spiritually: organicist worldview (Earth as living whole).

\textbf{Technocratic paradigm} (Francis): A worldview in which technical control and mastery are the dominant way of relating to reality; treats everything as a problem to be solved; excludes contemplative, spiritual, ethical ways.

\textbf{Spirituality}: Personal search for meaning, purpose, connectedness to a larger whole; three attributes: connectedness, self-transcendence, meaning. NOT necessarily tied to theism. (p.22--28)

\textbf{Magisterium} (Gould): Domain where a specific form of teaching authority holds the tools for meaningful discourse.

\textbf{Homeopathic dilution of religion} (Coyne): NOMA reduces religion to a version so weakened (no factual claims) it no longer resembles actual faith.

\textbf{Ichneumon wasps}: Parasitoid insects that eat hosts from inside; Darwin: cannot reconcile with benevolent God; evidence for amoral nature; used in Conflict model argument.
\end{multicols}

\vspace{8pt}

%=========================================================
\section{PAGINA-MARKERINGEN (POST-IT SYSTEEM)}
%=========================================================

\begin{infobox}
\textbf{Aanbevolen post-it labels voor je boek:}
\begin{itemize}[noitemsep]
  \item \textbf{``DEFS''} -- p.22--28: definities spiritualiteit, religie
  \item \textbf{``2.3.1''} -- p.51: relatie religie en spiritualiteit
  \item \textbf{``MISS CASE''} -- p.65: technologie en spiritualiteit
  \item \textbf{``BARBOUR''} -- p.76--80: 4 modellen (Conflict/Independence/Dialogue/Integration)
  \item \textbf{``NOMA''} -- NOMA sectie hfst 3: Gould's Non-Overlapping Magisteria
  \item \textbf{``WASP''} -- ichneumon wasps sectie: amoral nature, Darwin argument
  \item \textbf{``META AMB''} -- p.192: metaphysical ambiguity + Invisible Gardener
  \item \textbf{``GAIA''} -- hfst 7.1: Gaia hypothese + Doolittle kritiek
  \item \textbf{``EARTHRISE''} -- p.273--274: Earthrise foto + schrijf hier Blue Marble kritiek (Tom Hannes) + Axial/Post-axial uitleg \textbf{[EXAM juni 2025 Q1!]}
  \item \textbf{``TIKKUN''} -- hfst 6 judaïsme: Tikkun Olam + Bal Taschit
  \item \textbf{``KHALIFA''} -- hfst 6 islam: Khalifa + Tawhid + Akhrah
  \item \textbf{``LAUDATO''} -- p.254--262 + 320--327: Laudato Si' + moderate anthropocentrism
  \item \textbf{``BOEDDHISME''} -- p.265--266: Paticca-samuppada
  \item \textbf{``REN/WU''} -- p.267--269: Ren (Confucius) + Wuwei (Tao)
  \item \textbf{``ENVIRO RELIG''} -- hfst 8.2: environmentalisme als religie (Crichton)
  \item \textbf{``KV1 ANTW''} -- begin hfst 1: volledig antwoord sleutelvraag 1
  \item \textbf{``KV2a ANTW''} -- begin hfst 3: volledig antwoord sleutelvraag 2a
  \item \textbf{``KV2b ANTW''} -- begin hfst 7: volledig antwoord sleutelvraag 2b
  \item \textbf{``KV3 ANTW''} -- begin hfst 6: volledig antwoord sleutelvraag 3
  \item \textbf{``KV4 ANTW''} -- begin hfst 8: volledig antwoord sleutelvraag 4
  \item \textbf{``KV5 ANTW''} -- p.343--351: volledig antwoord sleutelvraag 5
  \item \textbf{``REVIEW''} -- achterste pagina: template review question
\end{itemize}
\end{infobox}

\end{document}
